\section{\texorpdfstring{هفته‌ی 10: مدل‌سازی بیزی 2: بررسی پیش‌بینی پسین و
\lr{Bayes
factor}}{هفته‌ی 10: مدل‌سازی بیزی 2: بررسی پیش‌بینی پسین و Bayes factor}}

\begin{quote}
\textbf{فصل 4: مدل‌سازی بیزی} (همراه با فصل 3: ارزیابی) · درس 1 ساعت ·
\lr{lab} 2 ساعت \textbf{داده‌ی اصلی:} داده‌های فراموشی هفته‌ی 2 (نمایی در
برابر توانی)، و یک شرکت‌کننده‌ی \lr{SDT} با شواهد ضعیف (آیا \(d' = 0\)
است؟).
\end{quote}

\begin{center}\rule{0.5\linewidth}{0.5pt}\end{center}

\subsection{واژه‌های کلیدی}

واژه‌های مهم این درس با معادل فارسی‌شان (حتی اگر در درس‌های قبل آمده
باشند).

{\def\LTcaptype{none} % do not increment counter
\begin{longtable}[]{@{}
  >{\raggedright\arraybackslash}p{(\linewidth - 4\tabcolsep) * \real{0.3333}}
  >{\raggedright\arraybackslash}p{(\linewidth - 4\tabcolsep) * \real{0.3333}}
  >{\raggedright\arraybackslash}p{(\linewidth - 4\tabcolsep) * \real{0.3333}}@{}}
\toprule\noalign{}
\begin{minipage}[b]{\linewidth}\raggedright
\lr{English}
\end{minipage} & \begin{minipage}[b]{\linewidth}\raggedright
فارسی
\end{minipage} & \begin{minipage}[b]{\linewidth}\raggedright
توضیح کوتاه
\end{minipage} \\
\midrule\noalign{}
\endhead
\bottomrule\noalign{}
\endlastfoot
\lr{posterior} & توزیع پسین \lr{(posterior)} & باور درباره‌ی پارامتر بعد
از دیدن داده \\
\lr{forgetting curve} & منحنی فراموشی & کاهش یادآوری با گذشت زمان \\
\lr{exponential} / \lr{power function} & تابع نمایی / توانی & دو شکل
رایج برای منحنی فراموشی و یادگیری \\
\lr{evidence} & شواهد & اطلاعاتی که برای تصمیم جمع می‌شود \\
\lr{sensitivity (d′)} & حساسیت & فاصله‌ی توزیع علامت و نویز \\
\lr{posterior predictive check (PPC)} & بررسی پیش‌بینی پسین & مقایسه‌ی
داده‌ی شبیه‌سازی‌شده از پسین با داده‌ی واقعی \\
\lr{marginal likelihood} / \lr{evidence} & درست‌نمایی حاشیه‌ای & احتمال
داده تحت مدل، میانگین روی پیشین \\
\lr{likelihood} & درست‌نمایی & احتمال داده به‌عنوان تابعی از پارامترها \\
\lr{Bayes factor (BF)} & عامل بیز \lr{(Bayes factor)} & نسبت
درست‌نمایی‌های حاشیه‌ای دو مدل \\
\lr{model flexibility} / \lr{complexity} & انعطاف‌پذیری / پیچیدگی مدل &
دامنه‌ی الگوهایی که مدل می‌تواند تولید کند \\
\lr{grid approximation} & تقریب شبکه‌ای \lr{(grid)} & محاسبه‌ی پسین روی
شبکه‌ای از مقادیر \\
\lr{Savage\lr{--}Dickey ratio} & نسبت \lr{Savage\lr{--}Dickey} & \lr{BF} از نسبت
چگالی پسین به پیشین در مقدار صفر \\
\lr{AIC} / \lr{BIC} & معیار اطلاعاتی آکائیکه / بیزی & برازش به‌علاوه‌ی
جریمه برای تعداد پارامتر \\
\lr{prior} & توزیع پیشین \lr{(prior)} & باور درباره‌ی پارامتر پیش از دیدن
داده \\
\lr{estimate} / \lr{estimator} & برآورد / برآوردگر & مقدار به‌دست‌آمده
برای پارامتر / روش به دست آوردن آن \\
\lr{posterior odds} & بخت پسین \lr{(odds)} & نسبت احتمال پسین دو مدل \\
\lr{fit} / \lr{fitting} & برازش & تنظیم پارامترها تا مدل با داده جور
شود \\
\lr{parameter} & پارامتر & کمیت قابل تنظیم مدل که بین افراد یا شرایط فرق
می‌کند \\
\lr{simulation} & شبیه‌سازی & تولید داده‌ی مصنوعی با اجرای مدل \\
\lr{noise} & نویز & تغییرات تصادفی در شواهد یا در داده \\
\lr{adequacy check} & بررسی کفایت مدل & آیا مدل ویژگی‌های اصلی داده را
بازتولید می‌کند \\
\lr{nested models} & مدل‌های تودرتو \lr{(nested)} & یک مدل حالت خاص مدل
دیگر است \\
\lr{trial} & کوشش \lr{(trial)} & یک بار انجام تکلیف توسط شرکت‌کننده \\
\lr{accuracy} & دقت & نسبت پاسخ‌های درست \\
\end{longtable}
}

\begin{center}\rule{0.5\linewidth}{0.5pt}\end{center}

\subsection{اهداف یادگیری}

در پایان این هفته، دانشجویان می‌توانند:

\begin{enumerate}
\def\labelenumi{\arabic{enumi}.}
\tightlist
\item
  توزیع‌های پیش‌بینی پسین را \textbf{تولید کنند} و بررسی‌های پیش‌بینی پسین
  \lr{(PPC)} نموداری و عددی \textbf{انجام دهند}.
\item
  \lr{marginal likelihood} و \lr{Bayes factor} را \textbf{تعریف کنند} و
  \textbf{توضیح دهند} چرا \lr{marginal likelihood} انعطاف مدل را خودکار
  جریمه می‌کند.
\item
  \lr{Bayes factor} را با انتگرال‌گیری شبکه‌ای و با نسبت چگالی
  \lr{Savage\lr{--}Dickey} \textbf{حساب کنند} و با تقریب \lr{BIC}
  \textbf{مقایسه کنند}.
\item
  حساسیت \lr{Bayes factor} به پیشین را \textbf{ارزیابی کنند} و آن را با
  پایداری برآوردهای پسین \textbf{مقایسه کنند}.
\item
  \lr{Bayes factor}ها را روی یک مقیاس درجه‌بندی‌شده‌ی شواهد \textbf{تفسیر
  کنند} و با بخت‌های پیشین \textbf{ترکیب کنند} تا احتمال‌های پسین مدل‌ها به
  دست آید.
\end{enumerate}

\begin{center}\rule{0.5\linewidth}{0.5pt}\end{center}

\subsection{مقدمه}

هفته‌ی قبل پسین‌ها را \emph{درون} یک مدل به دست آوردیم. دو سؤال باقی مانده
است.

\textbf{1. آیا مدل اصلاً خوب است؟} یک پسین می‌تواند باریک و مطمئن باشد حتی
وقتی مدل به‌شدت غلط است: مدل نمایی با کمال میل برای داده‌های هفته‌ی 2 یک
پسین دقیق برای \(\lambda\) می‌دهد. جواب بیزی این است که از مدل برازش‌شده
بخواهیم \textbf{داده‌ی جدید تولید کند} و ببینیم شبیه داده‌ی واقعی است یا
نه. این یک \emph{بررسی پیش‌بینی پسین} \lr{(posterior predictive check)}
است؛ همتای بیزی نمودار باقی‌مانده‌ها و \(G^2\).

\textbf{2. داده‌ها از کدام مدل بیشتر پشتیبانی می‌کنند؟} در هفته‌ی 7 مدل‌ها
را با \lr{AIC} و \lr{BIC} مقایسه کردیم که تعداد پارامترها را جریمه
می‌کنند. استنباط بیزی جوابی داخلی دارد: حساب کنیم هر مدل \emph{پیش از}
برازش داده‌ها را چقدر خوب پیش‌بینی کرده بود، با میانگین‌گیری روی همه‌ی
مقادیر پارامتری که پیشینش اجازه می‌دهد. این \textbf{\lr{marginal
likelihood}} است، و نسبت دو \lr{marginal likelihood} \textbf{\lr{Bayes
factor}} است. یک مدل منعطف پیش‌بینی‌هایش را روی مجموعه‌داده‌های ممکن زیادی
پخش می‌کند، پس به هرکدام احتمال کمتری می‌دهد. بنابراین \lr{Bayes factor}
تیغ اوکام را خودکار اجرا می‌کند، و به انعطاف شکل تابعی پاسخ می‌دهد، نه فقط
به \(k\).

\lr{Bayes factor}ها برای آزمون فرض صفر هم استفاده می‌شوند («آیا حساسیت
صفر است؟»). برخلاف مقدار \(p\)، می‌توانند شواهد \textbf{به نفع} فرض صفر
را هم کمّی کنند.

\paragraph{این موضوع کجا به کار می‌رود؟}

\begin{itemize}
\tightlist
\item
  \textbf{شاهد برای «بدون اثر».} مطالعه‌های تکرار \lr{(replication)} اغلب
  می‌پرسند آیا یک اثر وجود ندارد. \(p\)\lr{-value} نمی‌تواند از فرض صفر
  پشتیبانی کند؛ \lr{Bayes factor} می‌تواند («داده تحت فرض بدون اثر 8
  برابر محتمل‌تر است»).
\item
  \textbf{مجله‌ها و نرم‌افزارهای روان‌شناسی:} \lr{Bayes factor}ها با
  ابزارهایی مثل \lr{JASP} و بسته‌ی \lr{BayesFactor} در \lr{R} به‌طور معمول
  گزارش می‌شوند؛ هم برای \lr{t-test}، \lr{ANOVA} و همبستگی و هم برای
  مدل‌های شناختی.
\item
  \textbf{بررسی مدل در هر تحلیل بیزی:} \lr{posterior predictive check}
  گامی استاندارد پیش از اعتماد به هر نتیجه است (مثلاً در کارآزمایی‌های
  بالینی که به روش بیزی تحلیل می‌شوند و در مقاله‌های مدل‌سازی شناختی).
\item
  \textbf{مقایسه‌ی نظریه‌هایی با انعطاف متفاوت:} وقتی دو مدل تعداد پارامتر
  یکسانی دارند ولی یکی می‌تواند الگوهای داده‌ی خیلی بیشتری بسازد،
  \lr{Bayes factor} یکی از معدود ابزارهایی است که این را می‌بیند.
\end{itemize}

\begin{center}\rule{0.5\linewidth}{0.5pt}\end{center}

\subsection{جزوه‌ی درس}

\paragraph{به زبان ساده}

این هفته دو سؤال جدا داریم. \textbf{(1) آیا مدل من به اندازه‌ی کافی خوب
است؟} بگذارید مدل برازش‌شده داده‌ی تازه \emph{تصور کند}: مقدارهای معقول
پارامتر را از \lr{posterior} بردارید، یک آزمایش تازه شبیه‌سازی کنید و این
کار را بارها تکرار کنید. اگر داده‌ی واقعی شبیه یک داده‌ی تصوری معمولی
باشد، مدل کافی است؛ اگر داده‌ی واقعی بیرون از ابر داده‌های تصوری باشد، مدل
چیزی را از قلم انداخته است. این یک \textbf{\lr{posterior predictive
check}} است. \textbf{(2) داده از کدام‌یک از دو مدل بیشتر پشتیبانی می‌کند؟}
از هر مدل بپرسید، \emph{پیش از برازش}، داده‌ای را که واقعاً به دست آوردیم
چقدر محتمل می‌دانست، با میانگین گرفتن روی همه‌ی مقدارهای پارامتری که
\lr{prior} آن اجازه می‌دهد. این عدد \textbf{\lr{marginal likelihood}} است
و نسبت آن برای دو مدل \textbf{\lr{Bayes factor}} است. مدلی که روی همه‌چیز
شرط می‌بندد (خیلی منعطف) باید شرط‌هایش را پخش کند، پس وقتی داده جایی
می‌افتد که مدل ساده‌تر روی آن شرط بسته بود، کمتر برنده می‌شود. این همان تیغ
اوکام است که در نظریه‌ی احتمال تعبیه شده. نکته‌ی مهم: \lr{Bayes factor} به
\lr{prior} بستگی دارد، پس باید گزارش کنید که با \lr{prior}های معقول چطور
تغییر می‌کند.

\textbf{زمان‌بندی پیشنهادی (60 دقیقه):} بررسی پیش‌بینی پسین 15 ·
\lr{marginal likelihood} و \lr{Bayes factor} 20 · محاسبه‌ی \lr{Bayes
factor} (شبکه، \lr{BIC}، \lr{Savage\lr{--}Dickey)} 15 · حساسیت به پیشین و
مقایسه‌ی روش‌ها 10.

\subsubsection{1. بررسی‌های پیش‌بینی پسین}

\textbf{توزیع پیش‌بینی پسین} داده‌ی جدید \(\mathbf{y}^{\text{\lr{rep}}}\) را که
مدل برازش‌شده انتظار دارد توصیف می‌کند:

\[
p(\mathbf{y}^{\text{\lr{rep}}} \mid \mathbf{y}) = \int p(\mathbf{y}^{\text{\lr{rep}}} \mid \boldsymbol\theta)\, p(\boldsymbol\theta \mid \mathbf{y})\, d\boldsymbol\theta.
\]

\textbf{دستور شبیه‌سازی:} برای \(s = 1, \dots, S\)،
\(\boldsymbol\theta^{(s)}\) را از پسین بکشید و یک مجموعه‌داده‌ی
\(\mathbf{y}^{\text{\lr{rep}}(s)}\) از مدل با طرح آزمایش واقعی شبیه‌سازی کنید.
این کار \emph{هم} عدم قطعیت پارامترها و \emph{هم} نویز نمونه‌گیری را
منتقل می‌کند.

\textbf{بررسی‌ها:} - \textbf{نموداری:} داده‌ی مشاهده‌شده را روی بازه‌های
پیش‌بینی (مثلاً 95\% مرکزی \(y^{\text{\lr{rep}}}\) در هر شرط) رسم کنید.
انحراف‌های نظام‌مند (نقاط مشاهده‌شده که مرتب بیرون نوارها یا در یک
سمتشان‌اند) نشان می‌دهند مدل کجا شکست می‌خورد. - \textbf{عددی:} یک آماره‌ی
آزمون \(T(\mathbf{y})\) انتخاب کنید که ویژگی مورد نظرتان را نشان دهد
(مثلاً یادآوری در طولانی‌ترین فاصله‌ها، تندی افت اولیه). \textbf{مقدار
\(p\) پیش‌بینی پسین} برابر است با
\(p_B = P\big(T(\mathbf{y}^{\text{\lr{rep}}}) \ge T(\mathbf{y}) \mid \mathbf{y}\big)\).
مقادیر نزدیک 0 یا 1 نشان می‌دهند مدل نمی‌تواند این ویژگی را بازتولید کند.

\lr{PPC}ها کفایت \textbf{مطلق} را می‌سنجند، یعنی اینکه مدل به اندازه‌ی
کافی خوب است یا نه، نه اینکه کدام از دو مدل بهتر است. گام همیشگی هر
تحلیل بیزی‌اند (هفته‌ی 13).

\paragraph{خواندن ریاضی}

\begin{itemize}
\tightlist
\item
  \(\mathbf{y}^{\text{\lr{rep}}}\): داده‌ی «تکرارشده»، یک تکرار تصوری از
  آزمایش.
\item
  \(\int p(\mathbf{y}^{\text{\lr{rep}}} \mid \boldsymbol\theta)\, p(\boldsymbol\theta \mid \mathbf{y})\, d\boldsymbol\theta\):
  پیش‌بینی‌های مدل را روی همه‌ی مقدارهای پارامتر میانگین بگیر، هرکدام با
  وزن معقول بودنش در \lr{posterior.} در عمل این انتگرال با دستور
  شبیه‌سازی دومرحله‌ای جایگزین می‌شود: \(\boldsymbol\theta\) را از
  \lr{posterior} بکش و بعد با آن داده شبیه‌سازی کن.
\item
  \(p_B = P(T(\mathbf{y}^{\text{\lr{rep}}}) \ge T(\mathbf{y}) \mid \mathbf{y})\):
  کسری از داده‌های تصوری که آماره‌ی خلاصه‌ی \(T\) آن‌ها دست‌کم به اندازه‌ی
  داده‌ی واقعی بزرگ است. اگر فقط 1\% داده‌های تصوری یادآوری‌ای به بالایی
  مقدار مشاهده‌شده در 30 روز نشان دهند (\(p_B = 0.01\))، مدل نمی‌تواند این
  ویژگی را تولید کند.
\end{itemize}

\subsubsection{\texorpdfstring{2. \lr{marginal likelihood} و \lr{Bayes
factor}}{2. marginal likelihood و Bayes factor}}

\textbf{\lr{marginal likelihood}} (شواهد) مدل \(\mathcal{M}\)
درست‌نمایی‌ای است که روی پیشین میانگین‌گیری شده است:

\[
p(\mathbf{y} \mid \mathcal{M}) = \int p(\mathbf{y} \mid \boldsymbol\theta, \mathcal{M})\, p(\boldsymbol\theta \mid \mathcal{M})\, d\boldsymbol\theta.
\]

این احتمالی است که مدل \emph{پیش از دیدن} داده‌ی مشاهده‌شده به آن داده
بود. قاعده‌ی بیز در سطح مدل‌ها می‌دهد:

\[
\underbrace{\frac{p(\mathcal{M}_1 \mid \mathbf{y})}{p(\mathcal{M}_2 \mid \mathbf{y})}}_{\text{\lr{posterior odds}}}
= \underbrace{\frac{p(\mathbf{y} \mid \mathcal{M}_1)}{p(\mathbf{y} \mid \mathcal{M}_2)}}_{\mathrm{BF}_{12}}
\times \underbrace{\frac{p(\mathcal{M}_1)}{p(\mathcal{M}_2)}}_{\text{\lr{prior odds}}}.
\]

\(\mathrm{BF}_{12}\) ضریبی است که داده‌ها با آن بخت \(\mathcal{M}_1\)
نسبت به \(\mathcal{M}_2\) را تغییر می‌دهند، و
\(\mathrm{BF}_{21} = 1/\mathrm{BF}_{12}\).

\textbf{تیغ خودکار اوکام.} هر دو مدل باید احتمال کل 1 را روی همه‌ی
مجموعه‌داده‌های ممکن پخش کنند. یک مدل ساده پیش‌بینی‌هایش را روی چند
مجموعه‌داده متمرکز می‌کند؛ یک مدل منعطف آن‌ها را نازک پخش می‌کند. اگر داده‌ها
همان‌جایی بیفتند که مدل ساده پیش‌بینی کرده بود، مدل ساده \lr{marginal
likelihood} بالاتری می‌گیرد، حتی اگر مدل منعطف هم بعد از تنظیم
پارامترهایش بتواند آن‌ها را خوب برازش کند. برخلاف \lr{AIC/BIC}، این جریمه
به \emph{اینکه} مدل چقدر منعطف است (پراکندگی پیش‌بینی پیشین) بستگی دارد،
نه فقط به \(k\).

\textbf{یک مقیاس رایج شواهد} (برگرفته از دسته‌بندی کلاسیک \lr{Jeffreys}،
آن‌طور که \lr{Lee} و \lr{Wagenmakers} به کار می‌برند):

{\def\LTcaptype{none} % do not increment counter
\begin{longtable}[]{@{}ll@{}}
\toprule\noalign{}
\(\mathrm{BF}_{12}\) & شواهد به نفع \(\mathcal{M}_1\) \\
\midrule\noalign{}
\endhead
\bottomrule\noalign{}
\endlastfoot
1\lr{--}3 & حکایتی \lr{(anecdotal)} \\
3\lr{--}10 & متوسط \\
10\lr{--}30 & قوی \\
30\lr{--}100 & خیلی قوی \\
\lr{\textgreater} 100 & فوق‌العاده قوی \\
\end{longtable}
}

(مقادیر کمتر از 1 به همین شکل برای \(\mathcal{M}_2\) خوانده می‌شوند.) این
برچسب‌ها کمک‌کننده‌اند، نه آستانه. خود عدد را گزارش کنید.

\paragraph{خواندن ریاضی}

\begin{itemize}
\tightlist
\item
  \(p(\mathbf{y} \mid \mathcal{M}) = \int p(\mathbf{y} \mid \boldsymbol\theta)\,p(\boldsymbol\theta)\,d\boldsymbol\theta\):
  درست‌نمایی که روی \lr{prior} \textbf{میانگین} گرفته شده، نه بیشینه.
  مقایسه کنید: \lr{MLE} \emph{بهترین} مقدار پارامتر را برمی‌دارد؛
  \lr{marginal likelihood} \emph{میانگین وزن‌دار با \lr{prior}} روی همه‌ی
  مقدارها را. مدلی که \lr{prior} آن مقدارهای بد زیادی دارد، میانگین
  پایینی می‌گیرد.
\item
  \textbf{بخت پسین = \lr{Bayes factor} × بخت پیشین.} بخت \lr{(odds)}
  نسبت احتمال‌هاست: بخت 3 یعنی 3 به 1، یعنی احتمال 0.75. اگر دو مدل پیش
  از داده به یک اندازه معقول بودند (بخت پیشین = 1) و
  \(\mathrm{BF}_{12} = 3\) باشد، بعد از داده
  \(P(\mathcal{M}_1 \mid \mathbf{y}) = 3/(1+3) = 0.75\).
\item
  \textbf{مثال حل‌شده در یک خط:} با 7 درست از 10، مدل شانس داده را با
  احتمال 0.117 و مدل «هر دقتی» با احتمال 0.091 پیش‌بینی کرده بودند، پس
  \(\mathrm{BF}_{01} = 1.29\): داده بخت‌ها را تقریباً تغییر نمی‌دهد.
\end{itemize}

\subsubsection{\texorpdfstring{3. محاسبه‌ی \lr{Bayes
factor}}{3. محاسبه‌ی Bayes factor}}

\textbf{(الف) انتگرال‌گیری شبکه‌ای} (1 تا 2 پارامتر):
\(p(\mathbf{y}\mid\mathcal{M}) \approx \sum_g p(\mathbf{y}\mid\boldsymbol\theta_g)\, p(\boldsymbol\theta_g)\, \Delta_g\)،
که \(\Delta_g\) مساحت هر خانه است. آن را در فضای لگاریتمی
\lr{(log-sum-exp)} انجام دهید و بررسی کنید که جرم پیشین داخل شبکه \ensuremath{\approx} 1
باشد.

\textbf{(ب) تقریب \lr{BIC.}} برای \(N\) بزرگ،
\(\ln p(\mathbf{y}\mid\mathcal{M}) \approx -\tfrac12\mathrm{BIC}\)، پس

\[
\mathrm{BF}_{12} \approx \exp\big(-\tfrac12(\mathrm{BIC}_1 - \mathrm{BIC}_2)\big).
\]

سریع ولی خام است: پیشین واقعی را نادیده می‌گیرد (به‌طور ضمنی یک پیشین
\lr{«unit-information»} فرض می‌کند) و انعطاف شکل تابعی را هم.

\textbf{(ج) نسبت چگالی \lr{Savage\lr{--}Dickey}} (مدل‌های \lr{nested).} اگر
\(\mathcal{M}_0\) در \(\mathcal{M}_1\) مقدار \(\phi = \phi_0\) را ثابت
کند و پیشین‌های بقیه‌ی پارامترها یکی باشند، آن‌وقت

\[
\mathrm{BF}_{01} = \frac{p(\phi = \phi_0 \mid \mathbf{y}, \mathcal{M}_1)}{p(\phi = \phi_0 \mid \mathcal{M}_1)}:
\]

چگالی پسین در مقدار صفر تقسیم بر چگالی پیشین در مقدار صفر، هر دو تحت مدل
بزرگ‌تر. اگر داده‌ها \(\phi_0\) را محتمل‌تر از قبل کنند، شواهد به نفع فرض
صفر است. فقط یک پسین از \(\mathcal{M}_1\) لازم داریم.

(روش‌های دیگر، مثل \lr{bridge sampling} و \lr{thermodynamic integration}،
\lr{marginal likelihood} مدل‌های پیچیده را از خروجی \lr{MCMC} حساب
می‌کنند. این‌ها خارج از محدوده‌ی این درس‌اند.)

\paragraph{خواندن ریاضی}

\textbf{\lr{Savage\lr{--}Dickey}}
\(\mathrm{BF}_{01} = p(\phi_0 \mid \mathbf{y}) / p(\phi_0)\): به ارتفاع
\lr{posterior} و \lr{prior} در مقدار صفر \(\phi_0\) نگاه کنید (مثلاً
\(\theta = 0.5\)). اگر \lr{posterior} آنجا \emph{بلندتر} از \lr{prior}
باشد، داده مقدار صفر را معقول‌تر کرده، پس شواهد به نفع فرض صفر است. در
مثال حل‌شده \lr{prior} (صاف) در 0.5 ارتفاع 1 دارد و \lr{posterior} یعنی
\(\mathrm{Beta}(8,4)\) ارتفاع 1.29، پس \(\mathrm{BF}_{01} = 1.29\). یک
تصویر از دو منحنی و یک خط عمودی در \(\phi_0\) تمام چیزی است که برای
توضیح این لازم دارید.

\subsubsection{4. حساسیت به پیشین: برآورد در برابر آزمون}

برای \textbf{برآورد}، یک پیشین معقول با جمع شدن داده کم‌اهمیت‌تر می‌شود.
برای \textbf{\lr{Bayes factor}}، پیشین پارامتر آزمون‌شده \emph{هرگز}
بی‌اهمیت نمی‌شود. پهن‌تر کردن پیشین تحت \(\mathcal{M}_1\) پیش‌بینی‌هایش را
نازک‌تر پخش می‌کند و \(p(\mathbf{y}\mid\mathcal{M}_1)\) را کم می‌کند، پس یک
پیشین دلخواه مبهم \lr{Bayes factor} را، هر داده‌ای که باشد، به نفع فرض
صفر می‌کند (پارادوکس \lr{Jeffreys\lr{--}Lindley).} پیامدها: - پیشین‌ها برای
مقایسه‌ی مدل باید با دقت انتخاب شوند، ترجیحاً منعکس‌کننده‌ی اندازه‌ی اثرهای
واقع‌بینانه. - همیشه یک \textbf{تحلیل حساسیت} گزارش کنید: \lr{Bayes
factor} روی دامنه‌ای از پهنای معقول پیشین.

\subsubsection{5. مقایسه‌ی روش‌های ارزیابی مدل}

{\def\LTcaptype{none} % do not increment counter
\begin{longtable}[]{@{}
  >{\raggedright\arraybackslash}p{(\linewidth - 8\tabcolsep) * \real{0.2000}}
  >{\raggedright\arraybackslash}p{(\linewidth - 8\tabcolsep) * \real{0.2000}}
  >{\raggedright\arraybackslash}p{(\linewidth - 8\tabcolsep) * \real{0.2000}}
  >{\raggedright\arraybackslash}p{(\linewidth - 8\tabcolsep) * \real{0.2000}}
  >{\raggedright\arraybackslash}p{(\linewidth - 8\tabcolsep) * \real{0.2000}}@{}}
\toprule\noalign{}
\begin{minipage}[b]{\linewidth}\raggedright
روش
\end{minipage} & \begin{minipage}[b]{\linewidth}\raggedright
سؤال
\end{minipage} & \begin{minipage}[b]{\linewidth}\raggedright
چه چیزی را جریمه می‌کند
\end{minipage} & \begin{minipage}[b]{\linewidth}\raggedright
پیشین لازم است؟
\end{minipage} & \begin{minipage}[b]{\linewidth}\raggedright
هفته
\end{minipage} \\
\midrule\noalign{}
\endhead
\bottomrule\noalign{}
\endlastfoot
\(G^2\)، باقی‌مانده‌ها، \lr{PPC} & آیا مدل کافی است؟ & \lr{---} & \lr{PPC:} بله
& 4، 10 \\
\lr{LRT} & آیا پارامتر اضافه برازش را بیش از شانس بهتر می‌کند؟ (فقط
\lr{nested)} & \(\Delta k\) از راه \(\chi^2\) & نه & 7 \\
\lr{AIC} / \lr{CV} & کدام مدل داده‌ی جدید را بهتر پیش‌بینی می‌کند؟ &
\lr{overfitting} تخمینی & نه & 6\lr{--}7 \\
\lr{BIC} & به‌طور تقریبی: کدام مدل شواهد بیشتری دارد؟ & \(k \ln N\) &
ضمنی & 7 \\
\lr{Bayes factor} & داده‌ها بخت بین مدل‌ها را چقدر تغییر می‌دهند؟ &
پراکندگی پیش‌بینی پیشین (انعطاف واقعی) & \textbf{بله، و نتیجه به آن بستگی
دارد} & 10 \\
\end{longtable}
}

\begin{quote}
\textbf{پیشرفته: \lr{Bayes factor} در برابر معیارهای پیش‌بینی.}
\lr{marginal likelihood} را می‌شود به شکل حاصل‌ضرب پیش‌بینی‌های ترتیبی نوشت،
\(p(\mathbf{y}) = p(y_1)\,p(y_2\mid y_1)\cdots p(y_N\mid y_{1:N-1})\).
پس امتیاز می‌دهد که مدل هر نقطه‌ی داده را از روی نقاط قبلی، \emph{با شروع
از پیشین}، چقدر خوب پیش‌بینی می‌کند. \lr{LOO-CV} پیش‌بینی‌ها را از روی
(تقریباً) همه‌ی داده‌های دیگر امتیاز می‌دهد. به همین دلیل \lr{Bayes factor}
به پیشین وابسته است و \lr{LOO} تقریباً نه، و به همین دلیل ممکن است روی
داده‌ی یکسان با هم مخالف باشند.
\end{quote}

\paragraph{برای یک مهندس}

\begin{itemize}
\tightlist
\item
  \textbf{بخت پسین = \lr{BF} × بخت پیشین} همان \textbf{آشکارساز \lr{MAP}
  (بیز)} بین دو فرضیه است که در آن \lr{Bayes factor} نقش نسبت درست‌نمایی
  را دارد.
\item
  تفاوت با \textbf{\lr{GLRT}} (هفته‌های 4 و 7): با پارامترهای مجهول،
  \lr{GLRT} درست‌نمایی را روی آن‌ها \emph{بیشینه} می‌کند؛ \lr{Bayes factor}
  با یک \lr{prior} روی آن‌ها \emph{انتگرال} (میانگین) می‌گیرد. بیشینه کردن
  به فرضیه‌های منعطف امتیاز می‌دهد؛ میانگین گرفتن آن‌ها را خودکار جریمه
  می‌کند.
\item
  \textbf{پارادوکس \lr{Jeffreys\lr{--}Lindley}} رنگ‌وبوی مهندسی دارد: اگر
  بگذارید فرضیه‌ی جایگزین دامنه‌های سیگنال به‌طور نامعقول بزرگ را مجاز
  بداند \lr{(prior} خیلی پهن)، تقریباً هر مشاهده‌ای با «فقط نویز» بهتر
  توضیح داده می‌شود.
\item
  \textbf{\lr{posterior predictive check}} شبیه اعتبارسنجی یک مدل کانال
  است: سیگنال‌های دریافتی را از مدل برازش‌شده شبیه‌سازی کنید و آماره‌هایشان
  (هیستوگرام، خودهمبستگی، \lr{eye diagram)} را با اندازه‌گیری‌ها مقایسه
  کنید.
\end{itemize}

\paragraph{برای مدرس}

\begin{itemize}
\tightlist
\item
  \textbf{تیغ اوکام را با داستان شرط‌بندی توضیح دهید.} دو پیش‌بینی‌کننده
  هرکدام 100 ژتون دارند که باید روی نتیجه‌های ممکن بگذارند. یکی ژتون‌ها را
  روی همه‌چیز پخش می‌کند؛ دیگری بیشتر ژتون‌ها را روی یک محدوده‌ی باریک
  می‌گذارد. اگر نتیجه در آن محدوده‌ی باریک بیفتد، پیش‌بینی‌کننده‌ی متمرکز
  بیشتر می‌برد. \lr{marginal likelihood} همان مبلغ برد است.
\item
  \textbf{حساسیت به \lr{prior} را زنده نشان دهید.} در کد \lr{lab} پهنای
  \lr{prior} تحت \(\mathcal{M}_1\) را تغییر دهید و ببینید \lr{BF} چطور
  جابه‌جا می‌شود. این از هر هشداری قانع‌کننده‌تر است.
\item
  \textbf{دو سؤال را از هم جدا نگه دارید.} روی تخته بنویسید «آیا مدل
  کافی است؟ \ensuremath{\leftarrow} \lr{PPC»} و «کدام مدل بهتر است؟ \ensuremath{\leftarrow} \lr{BF} / \lr{CV} /
  \lr{AIC».} دانشجوها اغلب از \lr{BF} برای ادعای \emph{خوب} بودن یک مدل
  استفاده می‌کنند؛ \lr{BF} فقط می‌گوید \emph{از دیگری بهتر} است.
\item
  \textbf{سؤال‌های رایج.}

  \begin{itemize}
  \tightlist
  \item
    \emph{\lr{«BF} یا \(p\)\lr{-value}؟»} به سؤال‌های متفاوتی جواب
    می‌دهند. \lr{BF} دو فرضیه‌ی مشخص را مقایسه می‌کند و می‌تواند از فرض صفر
    پشتیبانی کند؛ \(p\)\lr{-value} فقط میزان غافلگیری تحت فرض صفر را
    می‌سنجد.
  \item
    \emph{\lr{«BF} = 2.5 معنادار است؟»} آستانه‌ای وجود ندارد؛ شواهد
    حکایتی \lr{(anecdotal)} است. عدد را گزارش کنید.
  \item
    \emph{«چرا \lr{BF} و \lr{AIC} با هم توافق ندارند؟»} هدف‌های متفاوتی
    دارند (شواهد در برابر پیش‌بینی) و \lr{BF} به \lr{prior} بستگی دارد.
  \end{itemize}
\end{itemize}

\begin{center}\rule{0.5\linewidth}{0.5pt}\end{center}

\subsection{مثال حل‌شده}

\subsubsection{\texorpdfstring{بخش 1: آیا یادآوری در حد شانس است؟ یک
\lr{Bayes factor} با
دست}{بخش 1: آیا یادآوری در حد شانس است؟ یک Bayes factor با دست}}

یک شرکت‌کننده در \(n = 10\) \lr{trial} بین دو گزینه انتخاب می‌کند و در
\(y = 7\) \lr{trial} درست جواب می‌دهد. آیا عملکرد در حد شانس است؟

\begin{itemize}
\tightlist
\item
  \(\mathcal{M}_0\): \(\theta = 0.5\) (بدون پارامتر آزاد).
\item
  \(\mathcal{M}_1\): \(\theta \sim \mathrm{Beta}(1, 1)\) (هر دقتی از پیش
  به یک اندازه محتمل است).
\end{itemize}

\textbf{قدم 1.}
\(p(y\mid\mathcal{M}_0) = \binom{10}{7} 0.5^{10} = 120/1024 = 0.1172\).

\textbf{قدم 2.} تحت پیشین یکنواخت، هر شمارش \(0, \dots, n\) از پیش به یک
اندازه محتمل است، پس
\(p(y\mid\mathcal{M}_1) = \int_0^1 \binom{10}{7}\theta^7(1-\theta)^3\, d\theta = \frac{1}{n+1} = \frac{1}{11} = 0.0909\).

\textbf{قدم 3.} \(\mathrm{BF}_{01} = 0.1172 / 0.0909 = 1.29\). داده‌ها
کمی \emph{به نفع شانس}‌اند؛ شواهدی حکایتی در هر دو جهت.

\textbf{قدم 4 (بررسی \lr{Savage\lr{--}Dickey).}} پسین تحت \(\mathcal{M}_1\)
برابر \(\mathrm{Beta}(8, 4)\) است که چگالی‌اش در 0.5 برابر 1.289 است.
چگالی پیشین در 0.5 برابر 1 است. نسبت 1.29 است، دقیقاً همان قدم 3.

\textbf{داده‌ی بیشتر:}

{\def\LTcaptype{none} % do not increment counter
\begin{longtable}[]{@{}
  >{\raggedright\arraybackslash}p{(\linewidth - 6\tabcolsep) * \real{0.2500}}
  >{\raggedright\arraybackslash}p{(\linewidth - 6\tabcolsep) * \real{0.2500}}
  >{\raggedright\arraybackslash}p{(\linewidth - 6\tabcolsep) * \real{0.2500}}
  >{\raggedright\arraybackslash}p{(\linewidth - 6\tabcolsep) * \real{0.2500}}@{}}
\toprule\noalign{}
\begin{minipage}[b]{\linewidth}\raggedright
داده
\end{minipage} & \begin{minipage}[b]{\linewidth}\raggedright
\(p\) (آزمون دوجمله‌ای دوطرفه)
\end{minipage} & \begin{minipage}[b]{\linewidth}\raggedright
\(\mathrm{BF}_{01}\)
\end{minipage} & \begin{minipage}[b]{\linewidth}\raggedright
خوانش
\end{minipage} \\
\midrule\noalign{}
\endhead
\bottomrule\noalign{}
\endlastfoot
7/10 & 0.34 & 1.29 & حکایتی، در هر دو جهت \\
60/100 & \ensuremath{\approx} 0.06 & 1.10 & هنوز حکایتی: داده‌ها به‌سختی شانس را از «کمی دقت»
جدا می‌کنند \\
70/100 & \lr{\textless} 0.001 & 0.0023 (\(\mathrm{BF}_{10} \approx 430\)) &
شواهد فوق‌العاده قوی علیه شانس \\
\end{longtable}
}

\textbf{تفسیر.} با 60/100، مقدار \(p\) نزدیک .05 پژوهشگران را وسوسه
می‌کند ادعای اثر کنند، ولی \lr{Bayes factor} می‌گوید داده‌ها تحت هر دو
فرضیه تقریباً به یک اندازه محتمل‌اند. دلیل: فرضیه‌ی جایگزین با پیشین
یکنواخت «پیش‌بینی کرده بود» که دقت می‌تواند هر جایی بین 0 و 1 باشد. دقت
مشاهده‌شده‌ی 0.6 فقط کمی بهتر با آن پیش‌بینی می‌شود تا با دقیقاً 0.5. یک
پیشین باریک‌تر و واقع‌بینانه‌تر تحت \(\mathcal{M}_1\) (مثلاً دقت‌های بین 0.5
و 0.8) \lr{Bayes factor} متفاوتی می‌داد؛ به همین دلیل حساسیت به پیشین
باید گزارش شود.

\subsubsection{بخش 2: احتمال‌های پسین مدل‌ها}

برای 7/10 با بخت پیشین 1: بخت پسین \(= 1.29\)، پس
\(P(\mathcal{M}_0 \mid y) = 1.29 / 2.29 = 0.56\). اگر یک داور از پیش
باور داشته باشد که عملکرد در حد شانس 4 برابر کمتر از عملکرد بالاتر از
شانس محتمل است (بخت پیشین 1:4)، آن‌وقت بخت پسین
\(= 1.29 \times 0.25 = 0.32\)، پس \(P(\mathcal{M}_0 \mid y) = 0.24\).
\lr{Bayes factor} برای همه یکی است؛ احتمال‌های پسین مدل‌ها به بخت‌های پیشین
بستگی دارند.

\begin{center}\rule{0.5\linewidth}{0.5pt}\end{center}

\subsection{\texorpdfstring{\lr{lab} (2 ساعت)}{lab (2 ساعت)}}

\subsubsection{اهداف}

\begin{enumerate}
\def\labelenumi{\arabic{enumi}.}
\tightlist
\item
  اجرای بررسی‌های پیش‌بینی پسین نموداری و عددی برای مدل‌های فراموشی نمایی و
  توانی.
\item
  محاسبه‌ی درست‌نمایی‌های حاشیه‌ای و \lr{Bayes factor} توانی در برابر نمایی
  با انتگرال‌گیری شبکه‌ای؛ مقایسه با \lr{BIC.}
\item
  آزمون \(d' = 0\) برای یک شرکت‌کننده‌ی \lr{SDT} با شواهد ضعیف با نسبت
  \lr{Savage\lr{--}Dickey}، و وارسی آن با محاسبه‌ی مستقیم.
\item
  ترسیم حساسیت \lr{Bayes factor} به پیشین و مقایسه‌ی آن با برآورد پسین.
\end{enumerate}

\textbf{زمان‌بندی:} بخش \lr{A} 35 دقیقه · بخش \lr{B} 25 دقیقه · بخش
\lr{C} 30 دقیقه · بخش \lr{D} 20 دقیقه · بخش \lr{E} 10 دقیقه.

\subsubsection{آماده‌سازی}

\begin{quote}
\textbf{\lr{Code} · کد:} \lr{\nolinkurl{code/week10.py}}, \lr{cell} 1 ·
سلول 1
\end{quote}

\textbf{خروجی مورد انتظار:}

\Needspace{7\baselineskip}
\begin{latin}
\begin{Verbatim}
exp: posterior mean A = 0.655, rate = 0.049
pow: posterior mean A = 0.867, rate = 0.417
\end{Verbatim}
\end{latin}

\subsubsection{\texorpdfstring{بخش \lr{A:} بررسی‌های پیش‌بینی پسین (35
دقیقه)}{بخش A: بررسی‌های پیش‌بینی پسین (35 دقیقه)}}

\begin{quote}
\textbf{\lr{Code} · کد:} \lr{\nolinkurl{code/week10.py}}, \lr{cell} 2 ·
سلول 2
\end{quote}

\textbf{خروجی مورد انتظار:}

\Needspace{7\baselineskip}
\begin{latin}
\begin{Verbatim}
exp: T(observed) = 0.267, mean T(rep) = 0.240, posterior predictive p = 0.302
pow: T(observed) = 0.267, mean T(rep) = 0.243, posterior predictive p = 0.320
\end{Verbatim}
\end{latin}

\textbf{تفسیر:} اول به نمودارها نگاه کنید. برای مدل نمایی، یادآوری
مشاهده‌شده در کوتاه‌ترین فاصله (42/50) \emph{بالای} نوار 95\% و یادآوری در
7 روز (13/50) \emph{پایین} آن است. مدل نمایی نمی‌تواند اول تند افت کند و
بعد ثابت شود. نوار مدل توانی همه‌ی نقاط مشاهده‌شده را در بر دارد.

با این حال بررسی عددی با \lr{\nolinkurl{T}} = میانگین یادآوری در سه
فاصله‌ی طولانی‌تر برای \textbf{هر دو} مدل مقادیر \(p_B\) معمولی (حدود 0.3)
می‌دهد. عدم برازش مدل نمایی در \emph{شکل} است (فاصله‌های اولیه در برابر
میانی)، که این آماره به آن نگاه نمی‌کند. \textbf{یک \lr{PPC} فقط شکست‌هایی
را تشخیص می‌دهد که آماره‌ی آزمونش به آن‌ها حساس است.} آماره‌هایی انتخاب کنید
که ویژگی‌های متمایزکننده‌ی نظریه‌ها را هدف بگیرند؛ در تمرین 2 یک آماره‌ی
«افت اول» مدل نمایی را لو می‌دهد (\(p_B \approx 0.02\)).

\subsubsection{\texorpdfstring{بخش \lr{B: Bayes factor} با انتگرال‌گیری
شبکه‌ای (25
دقیقه)}{بخش B: Bayes factor با انتگرال‌گیری شبکه‌ای (25 دقیقه)}}

\begin{quote}
\textbf{\lr{Code} · کد:} \lr{\nolinkurl{code/week10.py}}, \lr{cell} 3 ·
سلول 3
\end{quote}

\textbf{خروجی مورد انتظار:}

\Needspace{10\baselineskip}
\begin{latin}
\begin{Verbatim}
exp: log p(y|M) = -34.938   (prior mass inside grid = 0.9996)
pow: log p(y|M) = -22.872   (prior mass inside grid = 0.9991)

BF(pow vs exp) = exp(12.07) = 1.74e+05
BIC approximation: BF(pow vs exp) = exp(9.98) = 2.16e+04
\end{Verbatim}
\end{latin}

\textbf{تفسیر:} \(\mathrm{BF} \approx 1.7 \times 10^5\) شواهد فوق‌العاده
قوی به نفع مدل توانی در برابر نمایی است. \lr{BIC} همان حکم را می‌دهد ولی
با عدد دیگری (\(2 \times 10^4\)). چرا \lr{Bayes factor} شبکه‌ای و تقریب
\lr{BIC} با هم فرق دارند؟ (راهنمایی: هرکدام از کدام پیشین استفاده می‌کند،
و آیا دو مدل تحت پیشین‌های ما به یک اندازه منعطف‌اند؟)

\subsubsection{\texorpdfstring{بخش \lr{C: Savage\lr{--}Dickey:} آیا حساسیت
این شرکت‌کننده صفر است؟ (30
دقیقه)}{بخش C: Savage\lr{--}Dickey: آیا حساسیت این شرکت‌کننده صفر است؟ (30 دقیقه)}}

یک شرکت‌کننده در یک تکلیف تشخیص سخت: 12 \lr{hit} از 20 \lr{trial} علامت،
8 \lr{false alarm} از 20 \lr{trial} نویز. \(\mathcal{M}_1\):
\(d' \sim \mathcal{N}(0, 1)\)، \(c \sim \mathcal{N}(0, 1)\).
\(\mathcal{M}_0\): \(d' = 0\)، \(c \sim \mathcal{N}(0, 1)\).

\begin{quote}
\textbf{\lr{Code} · کد:} \lr{\nolinkurl{code/week10.py}}, \lr{cell} 4 ·
سلول 4
\end{quote}

\textbf{خروجی مورد انتظار:}

\Needspace{7\baselineskip}
\begin{latin}
\begin{Verbatim}
BF01 Savage-Dickey = 1.328;  BF01 direct = 1.328;  posterior mean d' = 0.444
MLE d' = 0.507
\end{Verbatim}
\end{latin}

دو محاسبه با هم توافق دارند؛ بررسی مفیدی برای اینکه هر دو درست پیاده‌سازی
شده‌اند.

\subsubsection{\texorpdfstring{بخش \lr{D:} حساسیت به پیشین (20
دقیقه)}{بخش D: حساسیت به پیشین (20 دقیقه)}}

\begin{quote}
\textbf{\lr{Code} · کد:} \lr{\nolinkurl{code/week10.py}}, \lr{cell} 5 ·
سلول 5
\end{quote}

\textbf{خروجی مورد انتظار:}

\Needspace{12\baselineskip}
\begin{latin}
\begin{Verbatim}
 prior SD of d'   BF01  BF10  posterior mean d'
           0.25  0.938 1.066              0.144
           0.50  0.972 1.028              0.313
           1.00  1.328 0.753              0.444
           2.00  2.314 0.432              0.496
           5.00  5.541 0.180              0.513
          10.00 11.011 0.091              0.515
\end{Verbatim}
\end{latin}

\textbf{به این دقت کنید:} با پهن‌تر شدن پیشین \(d'\)، میانگین پسین \(d'\)
سریع ثابت می‌شود، ولی \(\mathrm{BF}_{01}\) مدام به نفع فرض صفر بزرگ‌تر
می‌شود. یک پیشین خیلی پهن حساسیت‌های بزرگی را «پیش‌بینی می‌کند» که اتفاق
نیفتاده‌اند، پس فرضیه‌ی جایگزین جریمه می‌شود. \textbf{برای برآورد، پیشین
کم‌رنگ می‌شود؛ برای آزمون، نه.}

\subsubsection{\texorpdfstring{بخش \lr{E:} تفسیر (10
دقیقه)}{بخش E: تفسیر (10 دقیقه)}}

4 تا 5 جمله بنویسید: آیا این شرکت‌کننده علامت را تشخیص می‌دهد؟ \lr{Bayes
factor} را با پیشین انتخابی‌تان گزارش کنید و پیشین را توجیه کنید، به
تحلیل حساسیت اشاره کنید، و بگویید چه داده‌ی اضافه‌ای جواب را روشن‌تر می‌کرد.

\subsubsection{تمرین‌ها (به ترتیب سختی)}

\textbf{تمرین 1 \ensuremath{\star}.} با انتگرال‌گیری عددی
(\lr{\nolinkurl{scipy.integrate.quad}}) بررسی کنید که
\(\int_0^1 \binom{10}{7}\theta^7(1-\theta)^3 d\theta = 1/11\).

\textbf{تمرین 2 \ensuremath{\star}.} در بخش \lr{A} از یک آماره‌ی آزمون دیگر استفاده کنید:
افت یادآوری از اولین به دومین فاصله. کدام مدل (اگر هست) در این بررسی
شکست می‌خورد؟ (انتظار: مدل نمایی، \(p_B \approx 0.02\)؛ مدل توانی قبول
می‌شود.)

\textbf{تمرین 3 \ensuremath{\star}\ensuremath{\star} (حساسیت \lr{Bayes factor} فراموشی به پیشین).} بخش
\lr{B} را با پیشین‌های نرخ پهن‌تر (مقیاس × 3) و باریک‌تر (مقیاس ÷ 3) تکرار
کنید. آیا نتیجه‌ی «توانی \ensuremath{\gg} نمایی» باقی می‌ماند؟ چرا اینجا، برخلاف بخش
\lr{D}، این‌قدر پایدار است؟

\textbf{تمرین 4 \ensuremath{\star}\ensuremath{\star} (یک جایگزین آگاهانه).} در بخش \lr{C} پیشین \(d'\) را
با یک نیم‌نرمال \(d' \sim \mathcal{N}^+(0, 1)\) (فقط مقادیر مثبت؛
نرمال‌سازی دوباره را فراموش نکنید) جایگزین کنید. \lr{Bayes factor
Savage\lr{--}Dickey} را حساب کنید. چرا یک پیشین جهت‌دار نتیجه را تغییر می‌دهد؟

\textbf{تمرین 5 \ensuremath{\star}\ensuremath{\star}\ensuremath{\star} (پیشرفته: \lr{Bayes factor} برای مطالعه‌ی تکرار).}
برای مطالعه‌ی موردی هفته‌ی 8 (\(\mathcal{M}_0\) در برابر
\(\mathcal{M}_1\))، پیشین‌هایی برای \(\delta_0\)، \(\delta_1\) و
\(\kappa\) انتخاب کنید، \(\mathrm{BF}_{10}\) را با نسبت
\lr{Savage\lr{--}Dickey} در \(\delta_1 = 0\) با یک شبکه‌ی سه‌بعدی (درشت نگهش
دارید) حساب کنید و با تقریب \lr{BIC} مقایسه کنید.

\subsubsection{\texorpdfstring{\lr{lab check-off}}{lab check-off}}

نشان دهید: (1) دو نمودار \lr{PPC} با مقادیر \(p_B\)، (2) \lr{Bayes
factor} شبکه‌ای توانی در برابر نمایی و تقریب \lr{BIC} آن، و (3) جدول
حساسیت به پیشین از بخش \lr{D} با یک جمله تفسیر.

\begin{center}\rule{0.5\linewidth}{0.5pt}\end{center}

\subsection{\texorpdfstring{ارتباط با هوش مصنوعی و
\lr{LLM}ها}{ارتباط با هوش مصنوعی و LLMها}}

\textbf{آیا \lr{LLM}ها «نظریه‌ی ذهن» \lr{(theory of mind)} دارند؟} بعضی
مطالعه‌ها گزارش می‌دهند که \lr{LLM}ها تکلیف‌های باور نادرست \lr{(false
belief)} را با موفقیت انجام می‌دهند، یعنی آزمون‌های کلاسیکی که می‌سنجند آیا
کسی می‌فهمد دیگران می‌توانند باورهای نادرست داشته باشند. بعضی دیگر نشان
می‌دهند که تغییرهای کوچک در داستان‌ها عملکرد را فرو می‌ریزد {[}بررسی شود:
\lr{Kosinski,} 2023/2024; \lr{Ullman,} 2023{]}. بسیاری از این آزمون‌ها دو
گزینه‌ی جواب دارند، پس سطح شانس 50\% است، و تعداد آیتم‌ها کم است.
\lr{Bayes factor} به اولین سؤال درست جواب می‌دهد: \emph{آیا داده اصلاً از
عملکرد بالاتر از شانس، در مقایسه با شانس، پشتیبانی می‌کند؟}

\begin{quote}
\textbf{\lr{Code} · کد:} \lr{\nolinkurl{code/week10.py}}, \lr{cell} 6 ·
سلول 6
\end{quote}

\textbf{خروجی مورد انتظار:}

\Needspace{7\baselineskip}
\begin{latin}
\begin{Verbatim}
26/40 correct (65%): BF10 = 1.16
130/200 correct (65%): BF10 = 785.26
\end{Verbatim}
\end{latin}

همان 65\% با 40 آیتم اصلاً شاهدی نیست، و با 200 آیتم شاهدی قاطع است.
عملکرد بالاتر از شانس هم فقط قدم اول است: بررسی پیش‌بینی پسین این هفته
می‌پرسد مدل \emph{چطور} موفق می‌شود و شکست می‌خورد. مثلاً آیا دقیقاً روی
داستان‌های تغییریافته شکست می‌خورد، جایی که یک انسان شکست نمی‌خورد؟

\textbf{امتحانش کنید (بدون نیاز به \lr{API).}} - \ensuremath{\star} با یک پیشین متمرکز
روی مقادیر بالاتر از شانس تکرار کنید، مثلاً \lr{Beta(5,} 2) محدودشده به
\(\theta > 0.5\) (چگالی را دوباره نرمال کنید). \(\mathrm{BF}_{10}\) برای
26/40 چطور تغییر می‌کند؟ - \ensuremath{\star}\ensuremath{\star} یک مدل «میان‌بُر» \lr{(shortcut)} شبیه‌سازی
کنید که 90\% داستان‌های استاندارد و 30\% داستان‌های تغییریافته را درست
جواب می‌دهد (نیمی از هرکدام). دقت کلی‌اش 60\% است. یک بررسی پیش‌بینی پسین
طراحی کنید که میان‌بُر را آشکار کند. - \ensuremath{\star}\ensuremath{\star} بحث: «بالاتر از شانس در تکلیف‌های
باور نادرست» ادعایی در سطح محاسباتی \lr{Marr} است یا الگوریتمی؟

\begin{center}\rule{0.5\linewidth}{0.5pt}\end{center}

\subsection{اشتباهات رایج و برداشت‌های نادرست}

\begin{enumerate}
\def\labelenumi{\arabic{enumi}.}
\tightlist
\item
  \textbf{«پسین باریک یعنی مدل درست است.»} دقت همان کفایت نیست. همیشه یک
  \lr{PPC} اجرا کنید.
\item
  \textbf{استفاده از بررسی پیش‌بینی پسین برای انتخاب بین مدل‌ها.}
  \lr{PPC}ها کفایت را بررسی می‌کنند. برای مقایسه از \lr{Bayes factor}،
  \lr{CV} یا معیارهای اطلاعاتی استفاده کنید.
\item
  \textbf{\lr{«Bayes factor} احتمال پسین مدل است.»} \lr{Bayes factor}
  \emph{تغییر} بخت است. احتمال‌های پسین به بخت‌های پیشین هم نیاز دارند.
\item
  \textbf{استفاده از پیشین‌های خیلی مبهم در آزمون‌های \lr{Bayes factor.}}
  بی‌حد و مرز \lr{Bayes factor} را به سمت فرض صفر سوگیر می‌کنند (بخش
  \lr{D).} از پیشین‌های واقع‌بینانه استفاده کنید و حساسیت را گزارش کنید.
\item
  \textbf{\lr{Savage\lr{--}Dickey} با پیشین‌های ناهمخوان.} پیشین‌های پارامترهای
  مزاحم \lr{(nuisance)} باید در هر دو مدل یکی باشند، و چگالی پیشین باید
  \emph{در مقدار صفر} تحت مدل بزرگ‌تر ارزیابی شود.
\item
  \textbf{محاسبه‌ی \lr{marginal likelihood} روی شبکه‌ای که جرم پیشین را
  می‌بُرد.} بررسی کنید که جرم پیشین داخل شبکه نزدیک 1 باشد.
\item
  \textbf{\lr{«BF} = 2.9» و \lr{«BF} = 3.1» را کیفی متفاوت دانستن.}
  مقیاس شواهد پیوسته است؛ برچسب‌ها کمک‌کننده‌اند، نه آستانه‌ی تصمیم.
\end{enumerate}

\begin{center}\rule{0.5\linewidth}{0.5pt}\end{center}

\subsection{خودآزمایی}

\textbf{1.} توزیع پیش‌بینی پسین چگونه شبیه‌سازی می‌شود و چه چیزهایی را
ترکیب می‌کند؟ \textbf{پاسخ} پارامترها را از پسین بکشید و بعد با آن
پارامترها (و طرح آزمایش واقعی) از مدل داده شبیه‌سازی کنید. عدم قطعیت
درباره‌ی پارامترها را با نویز نمونه‌گیری داده‌ی جدید ترکیب می‌کند.

\textbf{2.} \lr{marginal likelihood} را با کلمات تعریف کنید.
\textbf{پاسخ} احتمال (چگالی) داده‌ی مشاهده‌شده تحت مدل، میانگین‌گیری‌شده روی
همه‌ی مقادیر پارامتر با وزن پیشین: اینکه مدل پیش از برازش این داده‌ها را
چقدر خوب پیش‌بینی کرده بود.

\textbf{3.} چرا \lr{marginal likelihood} حتی بدون شمردن پارامترها مدل‌های
منعطف را جریمه می‌کند؟ \textbf{پاسخ} احتمال روی مجموعه‌داده‌های ممکن باید
جمعش 1 شود. یک مدل منعطف پیش‌بینی‌های پیشینش را روی مجموعه‌داده‌های زیادی
پخش می‌کند و بنابراین به هرکدام، از جمله داده‌ی مشاهده‌شده، احتمال کمتری
می‌دهد؛ مگر اینکه داده جایی بیفتد که فقط مدل منعطف پیش‌بینی‌اش می‌کند.

\textbf{4.} \(\mathrm{BF}_{10} = 6\) و بخت پیشین 1:3 (علیه
\(\mathcal{M}_1\)) است. بخت پسین و \(P(\mathcal{M}_1\mid\mathbf{y})\)
چقدر است؟ \textbf{پاسخ} بخت پسین \(= 6 \times 1/3 = 2\)، پس
\(P(\mathcal{M}_1\mid\mathbf{y}) = 2/3 \approx 0.67\).

\textbf{5.} نسبت \lr{Savage\lr{--}Dickey} و شرایط استفاده از آن را بیان کنید.
\textbf{پاسخ}
\(\mathrm{BF}_{01} = p(\phi_0\mid\mathbf{y},\mathcal{M}_1)/p(\phi_0\mid\mathcal{M}_1)\):
چگالی پسین تقسیم بر چگالی پیشین در مقدار صفر تحت مدل بزرگ‌تر. به مدل‌های
\lr{nested} (فرض صفر \(\phi = \phi_0\) را ثابت می‌کند) و پیشین‌های یکسان
برای بقیه‌ی پارامترها نیاز دارد.

\textbf{6.} چرا پهن کردن پیشین پارامتر آزمون‌شده، حتی با داده‌ی ثابت،
\(\mathrm{BF}_{01}\) را افزایش می‌دهد؟ \textbf{پاسخ} یک پیشین پهن‌تر
احتمال بیشتری به اثرهای بزرگ و احتمال کمتری به اندازه‌ی اثرهای کوچک
سازگار با داده می‌دهد، پس \lr{marginal likelihood} فرضیه‌ی جایگزین کم
می‌شود در حالی که درست‌نمایی فرض صفر تغییری نمی‌کند.

\textbf{7.} \lr{BIC} مقدار \(\mathrm{BF}_{12} \approx 20\) می‌دهد، در
حالی که محاسبه‌ی شبکه‌ای با پیشین‌های شما 3 می‌دهد. کدام را گزارش می‌کنید و
چرا؟ \textbf{پاسخ} \lr{Bayes factor} شبکه‌ای را، چون از پیشین‌های واقعی
استفاده می‌کند و انعطاف شکل تابعی را در نظر می‌گیرد؛ \lr{BIC} تقریبی برای
نمونه‌ی بزرگ با یک پیشین ضمنی است. به اختلاف اشاره کنید و، در صورت امکان،
یک تحلیل حساسیت به پیشین بیاورید.

\begin{center}\rule{0.5\linewidth}{0.5pt}\end{center}

\subsection{تکلیف}

\subsubsection{\texorpdfstring{\lr{A4:} یک تحلیل بیزی با شبکه (1.0 نمره
از 20؛ انتشار هفته‌ی 10، موعد ابتدای درس هفته‌ی
12)}{A4: یک تحلیل بیزی با شبکه (1.0 نمره از 20؛ انتشار هفته‌ی 10، موعد ابتدای درس هفته‌ی 12)}}

\textbf{داده:} مجموعه‌داده‌ی \lr{A2}
(\lr{\nolinkurl{rng = np.random.default_rng(2024)}}، 9 فاصله، \(n = 20\)
آیتم؛ هفته‌ی 4 را ببینید).

\textbf{کارها.}

\begin{enumerate}
\def\labelenumi{\arabic{enumi}.}
\tightlist
\item
  \textbf{پیشین‌ها (1.5 امتیاز).} پیشین‌هایی برای \((A, \lambda)\) مدل
  نمایی و \((A, d)\) مدل توانی انتخاب و توجیه کنید. شبیه‌سازی‌های پیش‌بینی
  پیشین منحنی‌های فراموشی را برای هر دو مدل نشان دهید و در صورت نیاز
  پیشین‌ها را اصلاح کنید.
\item
  \textbf{پسین (2 امتیاز).} پسین‌های شبکه‌ای را برای هر دو مدل حساب کنید.
  میانگین‌های پسین، انحراف معیارها و \lr{HDI}های 95\%، و همبستگی پسین دو
  پارامتر را گزارش کنید. پسین‌های مشترک را رسم کنید.
\item
  \textbf{بررسی‌های پیش‌بینی پسین (2 امتیاز).} \lr{PPC} نموداری برای هر دو
  مدل، به‌علاوه‌ی یک \lr{PPC} عددی با آماره‌ی آزمونی که خودتان انتخاب و
  توجیه می‌کنید.
\item
  \textbf{\lr{Bayes factor} (2 امتیاز).}
  \(\mathrm{BF}_{\text{\lr{pow,exp}}}\) را با انتگرال‌گیری شبکه‌ای حساب کنید
  (جرم پیشین را بررسی کنید)، با تقریب \lr{BIC} مقایسه کنید، و با دو
  تنظیم پیشین دیگر تکرار کنید (تحلیل حساسیت).
\item
  \textbf{مقایسه با \lr{A2} (1 امتیاز).} خلاصه‌های پسین را با \lr{MLE}های
  \lr{A2} و عدم قطعیتشان مقایسه کنید.
\item
  \textbf{تفسیر (1.5 امتیاز، حداکثر 300 کلمه).} درباره‌ی فراموشی در این
  داده‌ها چه نتیجه‌ای می‌گیرید و این نتیجه‌گیری چقدر در برابر پیشین پایدار
  است؟
\end{enumerate}

\textbf{روبریک:} امتیازها مطابق بالا؛ داخل هر کار: درستی 60\%، نمودارها
و جدول‌ها 20\%، تفسیر 20\%. برای تحویل با تأخیر، روزی 10\% کسر می‌شود
(حداکثر 3 روز).

\begin{center}\rule{0.5\linewidth}{0.5pt}\end{center}

\subsection{مطالعه‌ی بیشتر}

\begin{itemize}
\tightlist
\item
  \textbf{\lr{Lee} \lr{\&} \lr{Wagenmakers} (2014)، فصل 7.} مقایسه‌ی بیزی
  مدل‌ها: درست‌نمایی‌های حاشیه‌ای، \lr{Bayes factor} و روش
  \lr{Savage\lr{--}Dickey}، با مثال‌های حل‌شده. \emph{(منبع اصلی.)}
\item
  \textbf{\lr{Farrell} \lr{\&} \lr{Lewandowsky} (2018)، فصل 11.} مقایسه‌ی بیزی
  مدل‌ها با \lr{Bayes factor}، در بافت مدل‌های شناختی.
\item
  \textbf{\lr{Kruschke} (2014)، فصل 10.} مقایسه‌ی مدل‌ها و مدل‌سازی
  سلسله‌مراتبی؛ مقدمه‌ای ساده بر \lr{Bayes factor} و حساسیتش به پیشین.
\item
  \textbf{\lr{Martin, Kumar} \lr{\&} \lr{Lao} (2021)، فصل 2.} تحلیل اکتشافی
  مدل‌های بیزی، از جمله بررسی‌های پیش‌بینی پسین.
\end{itemize}

\emph{پیش‌نمایش هفته‌ی 11:} شبکه‌های ما برای دو پارامتر کار کردند. برای
پنج، ده یا صد پارامتر به راه هوشمندانه‌تری برای کاوش پسین نیاز داریم:
\lr{Markov chain Monte Carlo.} اولین نمونه‌گیر خود را از صفر می‌سازید.
